% Proposed additive insert for Chapter 4 of the ProveIt manuscript.
% Requires the manuscript's existing \Li macro and theorem environments.
% Full proofs, branch details, and certificates are in the companion article
% ``Complementary depth, classical ladders, and exact shuffle ranks''.
% All new labels have the prefix cd: to avoid existing label collisions.

\section{Classical ladders and complementary depth}
\label{cd:sec:ladders}
Here the one-variable abbreviation is
$\Li_{s_1,\ldots,s_d}(z)=\Li_{s_1,\ldots,s_d}(z,1,\ldots,1)$.
Put
\[
 D=\mathbb C\setminus[0,\infty),\qquad
 q=(1-z)^{-1},\qquad L=\log q=-\log(1-z),
\]
with principal branches. At $z=\mathrm i$ these are
$q=(1+\mathrm i)/2$ and $L=-\tfrac12\log2+\mathrm i\pi/4$.
At $z=e^{2\pi\mathrm i/3}$ they are
$q=\tfrac12+\mathrm i\sqrt3/6$ and
$L=-\tfrac12\log3+\mathrm i\pi/6$.

\begin{theorem}[One-zero classical reduction]\label{cd:thm:onezero}
For $r\ge1$, define $F_r(z)=\Li_{2,\{1\}^{r-1}}(z)$. Then
\[
 F_r(z)=\sum_{k=1}^{r+1}
 \frac{(-1)^kL^{r+1-k}}{(r+1-k)!}\Li_k(q)
 +(-1)^r\zeta(r+1)-\frac{L^{r+1}}{(r+1)!}.
\]
For every $a,b\ge0$,
\[
 \Li_{\{1\}^a,2,\{1\}^b}(z)
 =\sum_{j=0}^{a}(-1)^j\binom{b+j+1}{j}
 \frac{L^{a-j}}{(a-j)!}F_{b+j+1}(z).
\]
These are exact functional identities on $D$.
\end{theorem}
\begin{proof}
For $z<0$, integrate the repeated $dt/(1-t)$ letters first. This gives
\[
 F_r(z)=\frac1{r!}\int_0^z\frac{[-\log(1-t)]^r}{t}\,dt.
\]
Set $u=(1-t)^{-1}$ and integrate by parts, retaining the endpoint
$u=1$, to obtain the first formula. More generally, integrating the
$a$ outer $dt/(1-t)$ letters gives
\[
 \Li_{\{1\}^a,2,\{1\}^b}(z)=\frac1{a!(b+1)!}
 \int_0^z\frac{(L+\log(1-t))^a[-\log(1-t)]^{b+1}}{t}\,dt.
\]
Binomial expansion proves the second formula. Analytic continuation
from $z<0$ proves the identities on $D$.
\end{proof}

The terminal-$2$ case simplifies to
\begin{align}
 \Li_{\{1\}^a,2}(z)
 ={}&-\frac{L^{a+2}}{(a+2)!}
 -\sum_{m=2}^{a+1}\frac{(m-1)L^{a+2-m}}{(a+2-m)!}\zeta(m)
 \notag\\
 &-L\Li_{a+1}(q)+(a+1)\Li_{a+2}(q)-(a+1)\zeta(a+2).
 \label{cd:eq:terminal}
\end{align}
Its derivative with respect to $L$ is the preceding case, and its
value tends to zero at $z=0^-$, which also proves it directly.

In particular, the three previously numerical weight-four Gaussian
triple displays are exact. Before specializing, the stronger complex
formulas are
\begin{align}
\Li_{2,1,1}(z)
={}&\Li_4(q)-L\Li_3(q)+\tfrac12L^2\Li_2(q)
 -\tfrac16L^3\Li_1(q)-\zeta(4)-\tfrac1{24}L^4,\notag\\
\Li_{1,2,1}(z)
={}&-3\Li_4(q)+2L\Li_3(q)-\tfrac12L^2\Li_2(q)
 +L\zeta(3)+3\zeta(4)-\tfrac1{24}L^4,\notag\\
\Li_{1,1,2}(z)
={}&3\Li_4(q)-L\Li_3(q)-2L\zeta(3)
 -\tfrac12L^2\zeta(2)-3\zeta(4)-\tfrac1{24}L^4.
\label{cd:eq:triples}
\end{align}
Their imaginary parts at $z=\mathrm i$ are precisely the three
printed formulas. The companion article supplies the classical
half-point evaluations and the exact symbolic substitution certificate.

\begin{theorem}[Complementary-depth upper bound]\label{cd:thm:complement}
A one-variable MPL of weight $w$ and depth $d$ has a finite rational
representation by terms
\[
 L^j\Li_{\boldsymbol a}(q)\zeta(\boldsymbol b),\qquad
 j+|\boldsymbol a|+|\boldsymbol b|=w,\qquad
 \operatorname{depth}(\boldsymbol a)+
 \operatorname{depth}(\boldsymbol b)\le w-d.
\]
The nonempty MZV indices are convergent. Empty indices denote $1$;
$|\boldsymbol a|$ denotes the sum of indices, not their number.
\end{theorem}
\begin{proof}
The binary word has $w-d$ zero letters. Under $u=(1-t)^{-1}$,
\[
 \frac{dt}{t}=-\frac{du}{u}-\frac{du}{1-u},\qquad
 \frac{dt}{1-t}=\frac{du}{u}.
\]
Every transformed word ends in zero and contains at most $w-d$ ones.
Transport its basepoint from $1$ to $0$ by splitting the word and
reversing each endpoint suffix. The reversed nonempty suffix begins
in zero, so its value at $1$ is convergent. Shuffle removal of trailing
zeros produces powers of $\log q$ and convergent MZVs, preserving the
total number of ones. The detailed endpoint justification and exact
word formula are in the companion article.
\end{proof}
This is a depth bound with logarithmic coefficients in an enlarged
argument family. It is not an assertion of reduction inside the
unchanged Gaussian or Eisenstein root-of-unity basket.

\subsection{The formal rank of the triple shuffle system}
In the formal shuffle algebra of binary words ending in one, the
weight-$w$, depth-$d$ indecomposable dimension is
\[
 \frac1w\sum_{k\mid\gcd(w,d)}\mu(k)\binom{w/k}{d/k}.
\]
There are $\binom{w-1}{d-1}$ original coordinates; subtracting the
above dimension gives the rank of the homogeneous product subspace.
This follows from the polynomial Lyndon-word basis and the primitive
necklace count. At $(w,d)=(6,3)$ the rank is seven on ten coordinates,
with the three Lyndon representatives $(4,1,1),(3,2,1),(3,1,2)$.
The supplied seven-row matrix is exact. None of this proves numerical
independence after taking values, real parts, or imaginary parts.

\subsection{A correction at mixed points}
The stuffle involution swaps both indices and colors:
\[
 \Li_{a,b}(x,y)+\Li_{b,a}(y,x)
 =\Li_a(x)\Li_b(y)-\Li_{a+b}(xy).
\]
Thus its antisymmetric complement is
$\Li_{a,b}(x,y)-\Li_{b,a}(y,x)$. It does not in general vanish
when $a=b$ and $x\ne y$. In fact,
\[
 \Li_{1,1}(\mathrm i,-1)-\Li_{1,1}(-1,\mathrm i)
 =\frac{\pi^2}{24}-\frac{\log^2 2}{4}
 +\mathrm i\left(G-\frac\pi2\log2\right)\ne0.
\]
This follows by integrating $-\log(1+x)/(1-x)$ from $0$ to
$\mathrm i$, and using the symmetric stuffle sum. Its real part is
positive. The original symmetric stuffle identities remain valid.
